\documentclass[11pt,a4paper]{article}

\usepackage[T1]{fontenc}
\usepackage[utf8]{inputenc}
\usepackage[ngerman]{babel}
\usepackage{lmodern}
\usepackage{microtype}
\usepackage[a4paper,margin=2.5cm]{geometry}
\usepackage{amsmath,amssymb,amsthm}
\usepackage{mathtools}
\usepackage{hyperref}

\hypersetup{
  colorlinks=true,
  linkcolor=blue,
  citecolor=blue,
  urlcolor=blue,
  pdftitle={Zur Hoffbauerschen Zwölf-Linie der Zetafunktion},
  pdfauthor={Thomas Hoffbauer und ChatGPT},
  pdfsubject={Bernoulli-Zahlen, Zetafunktion, Nennerstruktur},
  pdfkeywords={Riemannsche Zetafunktion, Bernoulli-Zahlen, von Staudt-Clausen, negative Spezialwerte}
}

\title{\textbf{Zur Hoffbauerschen Zwölf-Linie der Zetafunktion}\\
\large Nennerarithmetik ungerader negativer Spezialwerte und eine Strukturvermutung}
\author{Thomas Hoffbauer \and ChatGPT}
\date{\today}

\theoremstyle{definition}
\newtheorem{definition}{Definition}[section]
\newtheorem{beispiel}[definition]{Beispiel}

\theoremstyle{plain}
\newtheorem{satz}[definition]{Satz}
\newtheorem{proposition}[definition]{Proposition}
\newtheorem{lemma}[definition]{Lemma}
\newtheorem{korollar}[definition]{Korollar}
\newtheorem{vermutung}[definition]{Vermutung}
\newtheorem{empirischebehauptung}[definition]{Empirische Behauptung}

\theoremstyle{remark}
\newtheorem{bemerkung}[definition]{Bemerkung}

\begin{document}
\maketitle

\begin{abstract}
Für die ungeraden negativen Spezialwerte der Riemannschen Zetafunktion
\[
\zeta(-(2m-1))=-\frac{B_{2m}}{2m}
\]
untersuchen wir die vollständig gekürzten Nenner. Numerisch tritt der Nenner $12$ auffallend häufig auf. Diese Teilmenge von Indizes wird hier als \emph{Hoffbauersche Zwölf-Linie} bezeichnet. Wir formulieren eine exakte Teilbarkeitscharakterisierung dieser Indizes in Termen der Bernoulli-Zähler, diskutieren eine beobachtete Kongruenzstruktur und schlagen ein Modell für sogenannte Fast-Zwölf-Fälle vor, bei denen der Endnenner die Form $12r$ besitzt.
\end{abstract}

\section{Einleitung}

Die negativen ganzzahligen Spezialwerte der Riemannschen Zetafunktion sind durch die Bernoulli-Zahlen gegeben:
\[
\zeta(-n)=-\frac{B_{n+1}}{n+1}.
\]
Für ungerade negative Stellen folgt daraus
\[
\zeta(-(2m-1))=-\frac{B_{2m}}{2m}
\qquad (m\ge 1).
\]
Diese Werte sind rational. Damit stellt sich neben der Größe des Zählerbetrags insbesondere die Frage nach der Struktur der vollständig gekürzten Nenner.

Der prototypische Fall
\[
\zeta(-1)=-\frac{1}{12}
\]
ist klassisch. Numerisch zeigt sich jedoch, dass der Endnenner $12$ nicht isoliert bleibt, sondern an einer Reihe weiterer Indizes erneut auftritt. Diese Beobachtung motiviert die folgende strukturierende Terminologie.

\section{Grundlagen und Terminologie}

Schreiben wir die Bernoulli-Zahl $B_{2m}$ in vollständig gekürzter Form als
\[
B_{2m}=\frac{A_{2m}}{D_{2m}},
\]
wobei
\[
A_{2m}\in\mathbb Z,\qquad D_{2m}\in\mathbb N,\qquad \gcd(A_{2m},D_{2m})=1,
\]
so erhält man
\[
\zeta(-(2m-1))=-\frac{A_{2m}}{(2m)D_{2m}}.
\]

Nach dem Satz von von Staudt--Clausen gilt
\[
D_{2m}=\prod_{p-1\mid 2m} p,
\]
wobei das Produkt über alle Primzahlen $p$ läuft, deren Vorgänger $p-1$ ein Teiler von $2m$ ist. Der zugehörige Rohnenner ist daher
\[
R_m:=(2m)D_{2m}.
\]

\begin{definition}
Sei
\[
\zeta(-(2m-1))=\frac{P_m}{Q_m}
\]
in vollständig gekürzter Form mit
\[
P_m\in\mathbb Z,\qquad Q_m\in\mathbb N,\qquad \gcd(P_m,Q_m)=1.
\]
Dann nennen wir:
\begin{itemize}
\item $m$ einen \emph{Zwölf-Index}, falls $Q_m=12$,
\item $m$ einen \emph{Fast-Zwölf-Index}, falls $Q_m=12r$ mit einem ganzzahligen $r>1$.
\end{itemize}
Die Menge
\[
\mathcal Z_{12}:=\{m\in\mathbb N:Q_m=12\}
\]
heiße \emph{Hoffbauersche Zwölf-Linie}.
\end{definition}

\begin{bemerkung}
Die hier eingeführte Bezeichnung ist als Arbeitsbegriff zu verstehen. Sie beschreibt eine numerisch beobachtete Nennerklasse und erhebt keinen Anspruch auf bereits etablierte Terminologie.
\end{bemerkung}

\section{Eine exakte Charakterisierung}

\begin{satz}
Mit den obigen Bezeichnungen gilt
\[
Q_m=12
\quad\Longleftrightarrow\quad
\frac{(2m)D_{2m}}{12}\mid A_{2m}.
\]
\end{satz}

\begin{proof}
Aus
\[
\zeta(-(2m-1))=-\frac{A_{2m}}{(2m)D_{2m}}
\]
folgt für den vollständig gekürzten Nenner
\[
Q_m=\frac{(2m)D_{2m}}{\gcd(A_{2m},(2m)D_{2m})}.
\]
Somit ist $Q_m=12$ äquivalent zu
\[
\gcd(A_{2m},(2m)D_{2m})=\frac{(2m)D_{2m}}{12},
\]
und dies wiederum genau dann erfüllt, wenn
\[
\frac{(2m)D_{2m}}{12}\mid A_{2m}.
\]
\end{proof}

\begin{bemerkung}
Der Satz zeigt, dass die Hoffbauersche Zwölf-Linie nicht durch eine bloße Nennerbeobachtung beschrieben wird, sondern durch eine präzise Teilbarkeitsbedingung am Bernoulli-Zähler.
\end{bemerkung}

\section{Empirische Kongruenzstruktur}

Die numerische Untersuchung der ersten hundert ungeraden negativen Zetawerte zeigt eine deutliche Vorselektion der Zwölf-Indizes.

\begin{empirischebehauptung}
Für alle numerisch untersuchten Zwölf-Indizes $m\le 100$ gilt
\[
m\equiv 1 \text{ oder } 5 \pmod 6.
\]
Äquivalent dazu liegen die zugehörigen Bernoulli-Indizes in den Klassen
\[
2m\equiv 2 \text{ oder } 10 \pmod{12}.
\]
\end{empirischebehauptung}

\begin{bemerkung}
Im untersuchten Bereich tritt die Zwölf-Linie also nur bei solchen Indizes auf, die ungerade und nicht durch $3$ teilbar sind.
\end{bemerkung}

\begin{vermutung}[Dominanz der Klasse \texorpdfstring{$m\equiv 1\pmod 6$}{m ≡ 1 mod 6}]
Innerhalb der beiden zulässigen Klassen
\[
m\equiv 1,5 \pmod 6
\]
tritt die Hoffbauersche Zwölf-Linie signifikant häufiger in der Klasse
\[
m\equiv 1\pmod 6
\]
als in der Klasse
\[
m\equiv 5\pmod 6
\]
auf.
\end{vermutung}

\section{Fast-Zwölf-Fälle als Restphänomen}

Fast-Zwölf-Fälle besitzen Endnenner der Form
\[
Q_m=12r.
\]
Sie sind deshalb arithmetisch besonders interessant, weil sie nur um einen Restfaktor vom echten Zwölf-Fall entfernt sind.

\begin{definition}
Sei
\[
R_m=(2m)D_{2m}=12K_m.
\]
Dann nennen wir
\[
r_m:=\frac{Q_m}{12}
\]
den \emph{Restfaktor} des Index $m$.
\end{definition}

\begin{bemerkung}
Der Restfaktor misst genau den Anteil von $K_m$, der durch den Bernoulli-Zähler $A_{2m}$ nicht absorbiert wird. Der Zwölf-Fall entspricht daher gerade $r_m=1$.
\end{bemerkung}

\begin{vermutung}[Fast-Zwölf-Prinzip]
Fast-Zwölf-Indizes entstehen typischerweise dadurch, dass im Rohnenner
\[
(2m)D_{2m}=12K_m
\]
nicht alle Primfaktoren von $K_m$ im Bernoulli-Zähler $A_{2m}$ erscheinen. In vielen numerischen Beispielen bleibt lediglich ein einzelner kleiner Primfaktor oder ein kleines Primprodukt bestehen.
\end{vermutung}

\section{Von-Staudt-induzierte Restfaktoren}

Sobald für eine Primzahl $p$ die Bedingung
\[
p-1\mid 2m
\]
gilt, tritt $p$ automatisch im Nenner $D_{2m}$ auf. Damit liefert die von-Staudt--Clausensche Struktur eine natürliche Quelle möglicher Restfaktoren.

\begin{bemerkung}
Numerisch treten insbesondere Faktoren wie $11$ und $23$ gehäuft als Restfaktoren auf. Heuristisch lässt sich dies dadurch erklären, dass sie über die Teilbarkeitsbedingungen
\[
10\mid 2m
\qquad\text{bzw.}\qquad
22\mid 2m
\]
in den Nenner eingespeist werden. Ob diese Faktoren im Endnenner verbleiben, hängt dann ausschließlich davon ab, ob sie auch den Bernoulli-Zähler $A_{2m}$ teilen.
\end{bemerkung}

\begin{vermutung}[Von-Staudt-Kaskade]
Die Fast-Zwölf-Fälle werden häufig durch eine Kaskade von Primzahlen kontrolliert, die über Bedingungen der Form
\[
p-1\mid 2m
\]
in den Nenner eintreten. Die beobachteten Restfaktoren sind daher oft keine zufälligen Störungen, sondern strukturell erzeugte, vom Bernoulli-Zähler nicht absorbierte von-Staudt-Faktoren.
\end{vermutung}

\section{Der elementare Urfall \texorpdfstring{$\zeta(-1)=-1/12$}{zeta(-1)=-1/12}}

Der Wert
\[
\zeta(-1)=-\frac{1}{12}
\]
markiert den elementarsten Punkt der Zwölf-Linie. Aus
\[
B_2=\frac16
\]
folgt unmittelbar
\[
\zeta(-1)=-\frac{B_2}{2}=-\frac{1}{12}.
\]
Hier entsteht der Nenner $12$ direkt aus der Bernoulli-Struktur selbst. Es bedarf keiner komplizierten zusätzlichen Kürzung. In diesem Sinn kann der Wert $\zeta(-1)=-1/12$ als elementare Urgestalt der Hoffbauerschen Zwölf-Linie gelesen werden.

\section{Schlussbemerkung und Ausblick}

Die Hoffbauersche Zwölf-Linie erscheint als eine arithmetisch dünne Zone maximaler oder nahezu maximaler Kürzung innerhalb der ungeraden negativen Zetawerte. Der Rohnenner
\[
(2m)D_{2m}
\]
trägt die gesamte von-Staudt--Clausensche Information, während der Bernoulli-Zähler darüber entscheidet, welche Zusatzfaktoren im Endnenner überleben.

Daraus ergeben sich mehrere weiterführende Fragen:
\begin{itemize}
\item Besitzt die Zwölf-Linie eine asymptotisch beschreibbare Dichte innerhalb geeigneter Kongruenzklassen?
\item Lässt sich die Dominanz der Klasse $m\equiv 1\pmod 6$ theoretisch erklären?
\item Welche Rolle spielen reguläre und irreguläre Primzahlen für die Fast-Zwölf-Fälle?
\item Lässt sich die Restfaktor-Landschaft systematisch aus den von-Staudt-Daten und der Teilbarkeit der Bernoulli-Zähler rekonstruieren?
\end{itemize}

Damit deutet sich an, dass die beobachtete Zwölf-Linie nicht nur ein numerisches Kuriosum, sondern ein kleiner, präziser Zugang zur Feinarithmetik der Bernoulli-Zahlen und der negativen Spezialwerte der Riemannschen Zetafunktion sein könnte.

\end{document}
